\documentclass[10pt,a4paper]{article}
\usepackage{amsmath,amssymb,amsthm,graphicx}
\usepackage{hyperref}
% Card-local minimal preamble for S7 cards cut from sources/markov-quantile-jep-2022-177/source0.tex
% (doi:10.5802/jep.177, JEP 9 (2022), CC BY 4.0). Author text untouched.
% Omitted and why: the publisher cedram class (unavailable; assembler supplies
% article) and the publisher date commands (\datereceived/\dateaccepted/
% \dateepreuves/\datepublished), per the unavailable-class precedent.
\usepackage{amsmath,amssymb,amsthm}
\usepackage[scr=rsfs,cal=euler]{mathalfa}
\usepackage{xfrac}
\multlinegap0pt
\renewcommand{\theenumi}{\alph{enumi}}
\renewcommand{\theenumii}{\roman{enumii}}
\newenvironment{enumeratei}
{\bgroup\def\theenumi{\roman{enumi}}\def\theenumii{\arabic{enumii}}\begin{enumerate}}
{\end{enumerate}\egroup}
\makeatletter
\let\p@enumii\relax
\makeatother
\newcommand{\ldot}{\mathbin{.}}
\newcommand{\lcdots}{\ldot\cdots\ldot}
\newcommand{\RedefinitSymbole}[1]{%
\expandafter\let\csname old\string#1\endcsname=#1
\let#1=\relax
\newcommand{#1}{\csname old\string#1\endcsname\,}%
}
\RedefinitSymbole{\forall} \RedefinitSymbole{\exists}
\newcommand\mto{\mathchoice{\longmapsto}{\mapsto}{\mapsto}{\mapsto}}
\newcommand{\psfrac}[2]{\sfrac{(#1)}{#2}}
\newcommand{\spfrac}[2]{\sfrac{#1}{(#2)}}
\usepackage[all,cmtip]{xy}\let\labelstyle\textstyle
\usepackage{stmaryrd}
\usepackage{bbm}
\usepackage{mathtools}
\theoremstyle{plain}
\newtheorem{themBJ}{Theorem}
\renewcommand*{\thethemBJ}{\Alph{themBJ}}
\newtheorem{them}{Theorem}[section]
\newtheorem{pro}[them]{Proposition}
\newtheorem{pronotation}[them]{Proposition/Notation}
\newtheorem{prodefinition}[them]{Proposition/Definition}
\newtheorem{lemnotation}[them]{Lemma/Notation}
\newtheorem{cor}[them]{Corollary}
\newtheorem{lem}[them]{Lemma}
\newtheorem*{claim}{Claim}
\theoremstyle{definition}
\newtheorem{vocabulary}[them]{Vocabulary}
\newtheorem{defi}[them]{Definition}
\newtheorem{convention}[them]{Convention}
\newtheorem{notation}[them]{Notation}
\newtheorem{remnotation}[them]{Remark/Notation}
\newtheorem{definotation}[them]{Definition/Notation}
\newtheorem{remdefi}[them]{Remark/Definition}
\newtheorem{reminder-notation}[them]{Reminder/Notation}
\newtheorem{rem}[them]{Remark}
\newtheorem{ex}[them]{Example}
\newtheorem{reminder}[them]{Reminder}
\newtheorem*{problem*}{Main Problem}
\newtheorem*{note}{Note}
\newtheorem*{fact}{Fact}
\newcommand{\dd}{\mathrm{d}}
\newcommand{\rec}{\mathfrak R}
\newcommand{\id}{\mathrm{Id}}
\newcommand{\join}{\mathrm{Joint}}
\newcommand{\como}{\mathfrak{Q}}
\newcommand{\mq}{\mathfrak{MQ}}
\newcommand{\lev}{\mathrm{Lev}}
\newcommand{\ML}{\mathrm{LK}}
\newcommand{\IT}{\mathrm{IK}}
\newcommand{\disp}{\mathfrak{Disp}}
\newcommand{\argth}{\operatorname{Argth}}
\newcommand{\vol}{\operatorname{Vol}}
\newcommand{\parti}{\operatorname{Part}}
\newcommand{\idk}{\operatorname{id}}
\newcommand{\espace}{E}
\newcommand{\subd}{R}
\newcommand{\Tt}{{\mathcal T}}
\newcommand{\XX}{{\mathcal X}}
\newcommand{\EE}{{\mathcal E}}
\newcommand{\A}{{\mathcal A}}
\newcommand{\Var}{{\mathcal L}}
\newcommand{\Id}{\operatorname{Id}}
\newcommand{\one}{\mathbbm{1}}
\let\mb\relax
\let\mn\relax
\newcommand{\trans}{{}^{\mathrm{t}}\mkern-1mu}
\newcommand{\eqdef}{\coloneqq}
\newcommand{\eqdefup}{\coloneqq}
\newcommand{\espcdot}{\mbox{$\,\cdot\,$}}
\let\eps\varepsilon
\newcommand{\Lg}{\lambda|_{[0,1]}}
\newcommand{\la}{\lambda}
\renewcommand{\P}{\mathbb{P}}
\newcommand{\p}{\mathcal{P}}
\newcommand{\m}{\mathcal{M}}
\newcommand{\ma}{\mathrm{Marg}}
\newcommand{\cC}{\mathcal{C}}\let\C\cC
\newcommand{\R}{\mathbb{R}}
\newcommand{\Q}{\mathbb{Q}}
\newcommand{\N}{\mathbb{N}}
\newcommand{\Z}{\mathbb{Z}}
\newcommand{\AND}{\quad\text{and}\quad}
\newcommand{\E}{\mathbb{E}}
\newcommand{\law}{\operatorname{Law}}
\newcommand{\var}{\operatorname{Var}}
\newcommand{\sto}{\mathrm{sto}}
\newcommand{\stosup}{\operatorname{stosup}}
\newcommand{\stoinf}{\operatorname{stoinf}}
\newcommand{\lcsup}{\operatorname{losup}}
\newcommand{\lcinf}{\operatorname{loinf}}
\newcommand{\leqp}{\preceq_+}
\newcommand{\leqc}{\preceq_{C}}
\newcommand{\leqcs}{\preceq_{C,\sto}}
\newcommand{\leqs}{\preceq_{\text{\rm sto}}}
\newcommand{\geqs}{\succeq_{\text{\rm sto}}}
\newcommand{\e}{\mathrm{e}}
\newcommand{\pr}{\operatorname{proj}}
\newcommand{\pdec}{\mathcal{P}^{\scriptscriptstyle\searrow}}
\newcommand{\mdec}{\mathcal{M}^{\scriptscriptstyle\searrow}}
\newcommand{\leb}{\LL_{[0,1]}}
\newcommand{\op}{\mathopen}
\newcommand{\clo}{\mathclose}
\newcommand{\leqlc}{\preceq_{\rm lo}}
\newcommand{\geqlc}{\succeq_{\rm lo}}
\newcommand{\rsemibracket}{\mathclose{\rfloor\hspace*{-.9ex}\lceil}}
\newcommand{\lsemibracket}{\mathclose{\rfloor\hspace*{-.9ex}\lceil}}
\newcommand{\semi}{\mathclose{\rfloor\hspace*{-.8ex}\lceil}}
\newlength{\arraycolsepsauvegardegenerale}
\setlength{\arraycolsepsauvegardegenerale}{\arraycolsep}
\newlength{\arraycolsepsauvegardetemporaire}
\newlength{\tabcolsepsauvegardegenerale}
\setlength{\tabcolsepsauvegardegenerale}{\tabcolsep}
\newlength{\tabcolsepsauvegardetemporaire}
\newenvironment{nnarray}%
{\setlength{\arraycolsepsauvegardetemporaire}{\arraycolsep}\setlength{\arraycolsep}{0cm}\begin{array}}
{\end{array}\setlength{\arraycolsep}{\arraycolsepsauvegardetemporaire}}
\newenvironment{nntabular}%
{\setlength{\tabcolsepsauvegardetemporaire}{\tabcolsep}\setlength{\tabcolsep}{0cm}\begin{tabular}}
{\end{tabular}\setlength{\tabcolsep}{\tabcolsepsauvegardetemporaire}}
\providecommand{\LL}{\mathbf{L}}
\usepackage{cleveref}
\crefname{originaltheorem}{original theorem}{original theorems}
\Crefname{originaltheorem}{Original theorem}{Original theorems}
\makeatletter
\newlabel{them:kellerer_final}{{2.21}{source}{}{sourcepointer-1}{}}
\newlabel{them:kellerer_final@cref}{{[originaltheorem][0][]2.21}{[1][0][]source}{}{}{}}
\makeatother
\crefname{originaltheorem}{original theorem}{original theorems}
\Crefname{originaltheorem}{Original theorem}{Original theorems}
\makeatletter
\newlabel{rem:plan_kellerer}{{2.22}{source}{}{sourcepointer-2}{}}
\newlabel{rem:plan_kellerer@cref}{{[originaltheorem][0][]2.22}{[1][0][]source}{}{}{}}
\makeatother
\crefname{originaltheorem}{original theorem}{original theorems}
\Crefname{originaltheorem}{Original theorem}{Original theorems}
\makeatletter
\newlabel{them:existence2}{{2.15}{source}{}{sourcepointer-3}{}}
\newlabel{them:existence2@cref}{{[originaltheorem][0][]2.15}{[1][0][]source}{}{}{}}
\makeatother
\crefname{originaltheorem}{original theorem}{original theorems}
\Crefname{originaltheorem}{Original theorem}{Original theorems}
\makeatletter
\newlabel{them:existence5}{{2.15}{source}{}{sourcepointer-4}{}}
\newlabel{them:existence5@cref}{{[originaltheorem][0][]2.15}{[1][0][]source}{}{}{}}
\makeatother
\crefname{originaltheorem}{original theorem}{original theorems}
\Crefname{originaltheorem}{Original theorem}{Original theorems}
\makeatletter
\newlabel{them:existence}{{2.15}{source}{}{sourcepointer-5}{}}
\newlabel{them:existence@cref}{{[originaltheorem][0][]2.15}{[1][0][]source}{}{}{}}
\makeatother
\crefname{originaltheorem}{original theorem}{original theorems}
\Crefname{originaltheorem}{Original theorem}{Original theorems}
\makeatletter
\newlabel{them:existence3}{{2.15}{source}{}{sourcepointer-6}{}}
\newlabel{them:existence3@cref}{{[originaltheorem][0][]2.15}{[1][0][]source}{}{}{}}
\makeatother
\crefname{originaltheorem}{original theorem}{original theorems}
\Crefname{originaltheorem}{Original theorem}{Original theorems}
\makeatletter
\newlabel{lem:transitions_croissantes_limite}{{2.24}{source}{}{sourcepointer-7}{}}
\newlabel{lem:transitions_croissantes_limite@cref}{{[originaltheorem][0][]2.24}{[1][0][]source}{}{}{}}
\makeatother
\crefname{originaltheorem}{original theorem}{original theorems}
\Crefname{originaltheorem}{Original theorem}{Original theorems}
\makeatletter
\newlabel{them:existence4}{{2.15}{source}{}{sourcepointer-8}{}}
\newlabel{them:existence4@cref}{{[originaltheorem][0][]2.15}{[1][0][]source}{}{}{}}
\makeatother
\crefname{originaltheorem}{original theorem}{original theorems}
\Crefname{originaltheorem}{Original theorem}{Original theorems}
\makeatletter
\newlabel{lem:closed}{{2.23}{source}{}{sourcepointer-9}{}}
\newlabel{lem:closed@cref}{{[originaltheorem][0][]2.23}{[1][0][]source}{}{}{}}
\makeatother
\crefname{originaltheorem}{original theorem}{original theorems}
\Crefname{originaltheorem}{Original theorem}{Original theorems}
\makeatletter
\newlabel{rem:otherNs}{{2.18}{source}{}{sourcepointer-10}{}}
\newlabel{rem:otherNs@cref}{{[originaltheorem][0][]2.18}{[1][0][]source}{}{}{}}
\makeatother
\crefname{originaltheorem}{original theorem}{original theorems}
\Crefname{originaltheorem}{Original theorem}{Original theorems}
\makeatletter
\newlabel{rem:otherNsa}{{2.18}{source}{}{sourcepointer-11}{}}
\newlabel{rem:otherNsa@cref}{{[originaltheorem][0][]2.18}{[1][0][]source}{}{}{}}
\makeatother
\crefname{originaltheorem}{original theorem}{original theorems}
\Crefname{originaltheorem}{Original theorem}{Original theorems}
\makeatletter
\newlabel{them:existence1}{{2.15}{source}{}{sourcepointer-12}{}}
\newlabel{them:existence1@cref}{{[originaltheorem][0][]2.15}{[1][0][]source}{}{}{}}
\makeatother
\crefname{originaltheorem}{original theorem}{original theorems}
\Crefname{originaltheorem}{Original theorem}{Original theorems}
\makeatletter
\newlabel{rem:quantile_up}{{3.24}{source}{}{sourcepointer-13}{}}
\newlabel{rem:quantile_up@cref}{{[originaltheorem][0][]3.24}{[1][0][]source}{}{}{}}
\makeatother
\title{If ( t) t is a family of probability measures on / , increasing for , or , there is a Mark...}
\author{Boubel, Charles, Juillet, Nicolas}
\date{doi:10.5802/jep.177; CC BY 4.0}
\begin{document}
\maketitle
\setcounter{section}{1}
\setcounter{them}{17}
\setcounter{equation}{4}
\setcounter{themBJ}{2}
\noindent\textit{Editorial note.} All delivered mathematical text is the author's own, copied from the frozen source sources/markov-quantile-jep-2022-177/source0.tex; no statement or proof step was rewritten, reordered or removed. References inside the span are handed to the assembler as navigation stubs; load-bearing status is recorded per stub.\par\medskip
\noindent\textit{Editorial scope.} Single themBJ card cut from doi:10.5802/jep.177 (The Markov-quantile process attached to a family of marginals); statement and proof are the author own text at source lines 440-441 and 868-882.\par\medskip
\ifdefined\hypertarget\hypertarget{sourcepointer-1}{}\fi
\noindent\textit{Source reference.} Original theorem 2.21: sources/markov-quantile-jep-2022-177/source0.tex, source lines 833--833. navigation stub only: the target lives outside the counted span in the same article; whether it is load-bearing is judged per occurrence\par\medskip
\ifdefined\hypertarget\hypertarget{sourcepointer-2}{}\fi
\noindent\textit{Source reference.} Original theorem 2.22: sources/markov-quantile-jep-2022-177/source0.tex, source lines 836--836. navigation stub only: the target lives outside the counted span in the same article; whether it is load-bearing is judged per occurrence\par\medskip
\ifdefined\hypertarget\hypertarget{sourcepointer-3}{}\fi
\noindent\textit{Source reference.} Original theorem 2.15: sources/markov-quantile-jep-2022-177/source0.tex, source lines 781--781. navigation stub only: the target lives outside the counted span in the same article; whether it is load-bearing is judged per occurrence\par\medskip
\ifdefined\hypertarget\hypertarget{sourcepointer-4}{}\fi
\noindent\textit{Source reference.} Original theorem 2.15: sources/markov-quantile-jep-2022-177/source0.tex, source lines 787--787. navigation stub only: the target lives outside the counted span in the same article; whether it is load-bearing is judged per occurrence\par\medskip
\ifdefined\hypertarget\hypertarget{sourcepointer-5}{}\fi
\noindent\textit{Source reference.} Original theorem 2.15: sources/markov-quantile-jep-2022-177/source0.tex, source lines 775--775. navigation stub only: the target lives outside the counted span in the same article; whether it is load-bearing is judged per occurrence\par\medskip
\ifdefined\hypertarget\hypertarget{sourcepointer-6}{}\fi
\noindent\textit{Source reference.} Original theorem 2.15: sources/markov-quantile-jep-2022-177/source0.tex, source lines 783--783. navigation stub only: the target lives outside the counted span in the same article; whether it is load-bearing is judged per occurrence\par\medskip
\ifdefined\hypertarget\hypertarget{sourcepointer-7}{}\fi
\noindent\textit{Source reference.} Original theorem 2.24: sources/markov-quantile-jep-2022-177/source0.tex, source lines 858--858. navigation stub only: the target lives outside the counted span in the same article; whether it is load-bearing is judged per occurrence\par\medskip
\ifdefined\hypertarget\hypertarget{sourcepointer-8}{}\fi
\noindent\textit{Source reference.} Original theorem 2.15: sources/markov-quantile-jep-2022-177/source0.tex, source lines 785--785. navigation stub only: the target lives outside the counted span in the same article; whether it is load-bearing is judged per occurrence\par\medskip
\ifdefined\hypertarget\hypertarget{sourcepointer-9}{}\fi
\noindent\textit{Source reference.} Original theorem 2.23: sources/markov-quantile-jep-2022-177/source0.tex, source lines 854--854. navigation stub only: the target lives outside the counted span in the same article; whether it is load-bearing is judged per occurrence\par\medskip
\ifdefined\hypertarget\hypertarget{sourcepointer-10}{}\fi
\noindent\textit{Source reference.} Original theorem 2.18: sources/markov-quantile-jep-2022-177/source0.tex, source lines 814--814. navigation stub only: the target lives outside the counted span in the same article; whether it is load-bearing is judged per occurrence\par\medskip
\ifdefined\hypertarget\hypertarget{sourcepointer-11}{}\fi
\noindent\textit{Source reference.} Original theorem 2.18: sources/markov-quantile-jep-2022-177/source0.tex, source lines 816--816. navigation stub only: the target lives outside the counted span in the same article; whether it is load-bearing is judged per occurrence\par\medskip
\ifdefined\hypertarget\hypertarget{sourcepointer-12}{}\fi
\noindent\textit{Source reference.} Original theorem 2.15: sources/markov-quantile-jep-2022-177/source0.tex, source lines 779--779. navigation stub only: the target lives outside the counted span in the same article; whether it is load-bearing is judged per occurrence\par\medskip
\ifdefined\hypertarget\hypertarget{sourcepointer-13}{}\fi
\noindent\textit{Source reference.} Original theorem 3.24: sources/markov-quantile-jep-2022-177/source0.tex, source lines 1173--1173. navigation stub only: the target lives outside the counted span in the same article; whether it is load-bearing is judged per occurrence\par\medskip
\begin{themBJ}\label{them:c}If $(\mu_t)_{t\in \R}$ is a family of probability measures on\/ $\R$, increasing for $\leqc$, $\leqcs$ or $\leqs$, there is a Markov measure $P\in\ma((\mu_t)_t)$ which respectively is a martingale, is a submartingale or has increasing couplings.
\end{themBJ}

\begin{proof}[Proof of Theorem \ref{them:c}] Take $(\mu_t)_t\in\p(\R)^\R$, increasing for $\leqc$ (case \eqref{pf:thCa}), $\leqcs$ (case~\eqref{pf:thCb}) or for $\leqs$ (case \eqref{pf:thCc}) to prove the corresponding cases of Theorem \ref{them:c}. In the sketch of proof of Theorem \ref{them:kellerer_final} given in Remark \ref{rem:plan_kellerer}, replace $\mathcal{N}^\ML_{s,t}$ by $\mathcal{N}^\IT_{s,t}$ and introduce, similarly as defined in Remark \ref{rem:plan_kellerer} for cases \eqref{pf:thCa} and \eqref{pf:thCb}, the spaces~$\mathcal{N}_{s,t}'^\IT$---equal to
$\{P\in\mathcal{N}^\IT_{s,t}:\,P(\{(x,y)\in\R^2:\,x\leqslant y\})=1\}$ in case \eqref{pf:thCc}.

Properties \eqref{them:existence2}--\eqref{them:existence5} of Theorem \ref{them:existence} are satisfied by $(\mathcal{N}^\IT_{s,t})_{s<t}$: \eqref{them:existence2} by definition, \eqref{them:existence3}~by Lemma \ref{lem:transitions_croissantes_limite}, \eqref{them:existence4} is straightforward and \eqref{them:existence5} by Lemma \ref{lem:closed}. By Remark \ref{rem:otherNs}\eqref{rem:otherNsa} we are left with showing \eqref{them:existence1}, \eqref{them:existence3} and \eqref{them:existence4} for $(\mathcal{N}_{s,t}'^\IT)_{s<t}$. Plainly, the conditions defining the $\mathcal{N}_{s,t}'^\IT$ as subspaces are closed and stable by composition, \eqref{them:existence3} and \eqref{them:existence4} follow. For~\eqref{them:existence1}, in our three cases:
\begin{enumerate}
\item\label{pf:thCa}
By \cite[\S 3.1]{BHS}, if $P\in\ma(\mu_s,\mu_t)$ is a martingale transport plan, then $P\in\mathcal{N}_{s,t}^\ML\Leftrightarrow P\in\mathcal{N}_{s,t}^\IT$. Therefore $\mathcal{N}_{s,t}'^\IT\neq\varnothing$ if and only if $\mathcal{N}_{s,t}'^\ML\neq\varnothing$, which Kellerer proved.

\item\label{pf:thCb}
The element of $\mathcal{N}_{s,t}'^\ML$ Kellerer built in \cite[Def.\,7 \& Th.\,2]{Ke72} is in $\mathcal{N}_{s,t}'^\IT$ when $\mu_s\leqc\mu_t$.

\item\label{pf:thCc}
As explained in Remark \ref{rem:quantile_up}, $\como(\mu_s,\mu_t)\in \mathcal{N}_{s,t}'^\IT$.\qedhere
\end{enumerate}
\end{proof}

% Retained bibliography: cited entries copied verbatim from the
% publisher bibliography of 10.5802/jep.177.
\begin{thebibliography}{99}
\bibitem{BHS} M. Beiglböck, M. Huesmann \& F. Stebegg - “Root to Kellerer”, in Séminaire de Probabilités XLVIII, Lect. Notes in Math., vol. 2168, Springer, Cham, 2016, p. 1-12
\bibitem{Ke72} H. G. Kellerer - “Markov-Komposition und eine Anwendung auf Martingale”, Math. Ann. 198 (1972), p. 99-122
\end{thebibliography}
\end{document}
